\documentclass[12pt]{article}

%--- Packages ---
\usepackage[margin=1in]{geometry}
\usepackage{amsmath, amssymb}
\usepackage{setspace}
\usepackage{natbib}
\usepackage{hyperref}
\usepackage{booktabs}
\usepackage{graphicx}
\usepackage{caption}

\doublespacing

\title{\textbf{A Quantity Theory of Money Approach to Cryptocurrency Valuation:\\
Store of Value, Medium of Exchange, and the Economic Moat Signal}}

\author{Moazzam Khoja\\
\small University of Texas at San Antonio\\
\small \texttt{moazzam.khoja@utsa.edu}}

\date{June 2026\\[0.5em]
\small \textit{Preliminary. Comments welcome.}\\
\small \textit{Generated with AI assistance under AFA 2027 Special Session guidelines.}}

\begin{document}

\maketitle

\begin{abstract}
We apply the Quantity Theory of Money (QTM) to the cross-sectional valuation of
cryptocurrencies and protocol tokens. Treating each digital asset as the currency of a
self-contained economy, we decompose market capitalization into a Medium of Exchange
(MoE) component, priced by the free-circulating velocity of supply, and a Store of Value
(SoV) residual. The key innovation is $\lambda$---the fraction of total supply that is
locked or staked---which partitions the circulating supply into economically active and
inert components. A central result is that the SoV/MoE ratio simplifies to
$\lambda/(1-\lambda)$, directly observable from on-chain data without velocity estimation.
For native blockchain coins, high $\lambda$ reflects a convenience yield and seigniorage
signal grounded in monetary search theory. For governance and utility tokens, which lack
a medium-of-exchange function, $\lambda$ itself is an economic moat signal, capturing
token holder fundamental conviction in platform growth. We show that the conviction
signal alone positively predicts future cross-sectional returns, and that this
predictability is substantially stronger when the asset is cheap relative to its
fundamental economic activity---measured by our Growth-Levelized NVT ratio for coins and
the Growth-Levelized NV/TVL ratio for tokens---supporting the interpretation that valuation identifies
when the conviction signal is not yet priced in. A long Stars/short Avoid portfolio
generates positive factor-adjusted alpha after controlling for market, size, and momentum
factors.
\end{abstract}

\newpage

%=============================================================
\section{Introduction}
\label{sec:intro}
%=============================================================

What determines the value of a cryptocurrency? The dominant approach in the literature
has been to adapt discounted cash flow (DCF) models, searching for a dividend-like
flow---transaction fees, staking rewards, governance rights---that can be capitalized
into a present value \citep{liu2021risks, pagnotta2022equilibrium, biais2019blockchain}.
This approach, however, mischaracterizes the fundamental nature of digital assets.
Cryptocurrencies and protocol tokens exhibit currency-like behavior: each coin or token
is the monetary instrument of a self-contained digital economy, and its value is governed
by the same forces that determine the value of any currency---the scale of economic
activity it enables and the efficiency with which it circulates. The correct framework,
we argue, is the Quantity Theory of Money \citep{fisher1911}.

The Fisher equation, $MV = PQ$, relates the stock of money ($M$) and its velocity ($V$)
to the nominal value of economic output ($PQ$) \citep{fisher1911}. Applied to digital
assets, a cryptocurrency's market capitalization proxies its money supply, and the
economic output it supports---on-chain transaction volume for native coins, protocol
throughput for tokens---is its $PQ$. This is not merely an analogy.
\citet{cong2021tokenomics} develop a formal dynamic model in which token prices aggregate
heterogeneous users' transactional demand, with equilibrium prices governed by platform
adoption---a formulation directly consistent with the QTM framework we develop here.
\citet{sockin2020model} similarly model tokens as platform membership instruments,
showing that the token price reflects the present value of the convenience yield accruing
to participants: the non-pecuniary benefit of network membership, including access to
platform services, network liquidity, and future transactional capability.

Before proceeding, we define the two components of value our framework decomposes. The
\textit{Medium of Exchange} (MoE) value of a digital asset is the value attributable to
its current transactional function: the economic activity it enables today, priced by the
velocity with which it circulates. The \textit{Store of Value} (SoV) value is the premium
agents are willing to pay today by forgoing current exchange utility---the transactions,
goods, services, or alternative investments they could obtain by spending the asset
now---in anticipation of future value appreciation. This appreciation is expected to arise
from network growth, expanding user adoption, increasing utility, and innovation. An agent
holding Ethereum rather than spending it is implicitly paying an intertemporal price: she
sacrifices today's exchange utility for her conviction that the network's future worth
exceeds today's use value. The SoV component is the aggregate market valuation of this
forward-looking conviction.

The QTM framework requires a critical extension to operationalize this decomposition. Not
all of a cryptocurrency's circulating supply is economically active. A substantial and
growing fraction is locked in staking contracts, pledged as governance collateral, or held
in long-term custody with near-zero velocity. These holdings are economically inert from
a transactional standpoint: they contribute to the money supply but perform no current
exchange function. Conflating active and inactive supply distorts velocity estimates and,
consequently, misprices the economic content of market capitalization. We address this by
introducing $\lambda$, defined as the proportion of the total supply that is locked or
staked---including tokens committed to network validation, governance contracts, and other
long-term holding functions. The formal derivations are provided in
Appendix~\ref{app:derivations}. The key result is that the SoV/MoE ratio simplifies to
\begin{equation}
\frac{\text{MC}_{\text{SoV}}}{\text{MC}_{\text{MoE}}} = \frac{\lambda}{1-\lambda},
\label{eq:sov_moe}
\end{equation}
a monotone, directly observable function of $\lambda$ alone, requiring no velocity
estimation.

Why does $\lambda$ matter beyond its role in the decomposition? For native blockchain
coins---Ethereum, Solana, and their peers---high $\lambda$ is grounded in three
complementary theoretical traditions. First, \textit{monetary search theory}
\citep{kiyotaki1993search, lagos2005unified} establishes that money emerges as a
coordination equilibrium: an asset becomes valuable because rational agents expect others
to hold and accept it. High $\lambda$ is the blockchain-observable, publicly verifiable
signal of this coordination commitment. Second, the \textit{convenience yield}
\citep{sockin2020model, krishnamurthy2012aggregate} is the non-pecuniary benefit agents
derive from holding a monetary asset: access to network security, platform liquidity, and
future transactional capability. Just as holders of U.S. Treasury securities accept lower
yields in exchange for liquidity services \citep{nagel2016liquidity}, holders of dominant
blockchain coins accept lower required returns because network membership has value beyond
current transactions. Crucially, the SoV premium is the present value of expected
\textit{future} convenience yields: agents who lock coins today are paying for anticipated
growth in the network's membership value. Third, the \textit{seigniorage} channel
\citep{cong2024staking} anchors this signal in rational optimization: in proof-of-stake
networks, newly minted tokens distributed to stakers constitute monetary seigniorage, and
the staking decision reveals belief that this income plus anticipated price appreciation
exceeds the opportunity cost of locking capital. These three channels are complementary:
coordination theory explains why high $\lambda$ is self-reinforcing, convenience yield
provides the holding motive, and seigniorage grounds the rational optimization.

Governance and utility tokens present a related but distinct case, and the distinction
matters for how the framework is applied. A governance token such as UNI is the monetary
instrument of the Uniswap protocol economy, but it has no medium-of-exchange function:
the token does not need to change hands to facilitate Uniswap trades, so there is no MoE
component to separate from an SoV premium. With the MoE denominator absent, the SoV/MoE
ratio is undefined for tokens, and we use $\lambda$ itself to convey the same conviction
content directly. The protocol's economic output is still well defined---decentralized
exchange volume, liquidity deployed, fees generated---and whether protocol revenue
currently accrues to token holders is immaterial, just as it is immaterial whether U.S.
GDP flows to dollar holders. This framing aligns with \citet{cong2021tokenomics}, who
price tokens on the present value of future platform utility and adoption, and
\citet{schilling2019bitcoin}, who model the value of digital currencies as functions of
the economic activity they support. For governance tokens, $\lambda$ carries a different
signal. Staking a governance token into a DAO contract renders it illiquid: staked tokens
cannot be sold while locked. An agent who voluntarily incurs this illiquidity cost reveals
a preference for future appreciation over present liquidity---a choice grounded in
assessments of network effects, utility expansion, competitive positioning, and innovation
trajectory. Drawing on \citet{hirschman1970exit}, holders who exercise \textit{voice}
(governance participation) rather than \textit{exit} (selling) reveal a preference for
the protocol's future that passive buyers do not. \citet{edmans2011governance} formalize
this in a corporate setting: blockholders who engage rather than exit create governance
value precisely because engagement is costly and therefore credible. Applied to tokens,
active governance participation is a costly signal \citep{spence1973job} of fundamental
conviction in platform growth, making high participation rates an economic moat indicator.

We operationalize these ideas in two metrics per asset type. For coins, the
\textit{SoV/MoE ratio}---equal to $\lambda/(1-\lambda)$---measures the balance between
forward-looking conviction and current transactional utility, and is our primary return
predictor; for tokens, $\lambda$ itself plays this role, as described above. The second
metric normalizes market capitalization by fundamental economic activity, and its
denominator likewise adapts to the asset type. For coins, the \textit{Growth-Levelized
NVT ratio} normalizes market capitalization by the present value of the network's
expected transaction stream $PQ$ (constructed analogously to a PMT-based discounted cash
flow, as detailed in Section~\ref{sec:measures}). For tokens, protocol throughput is
measured by Total Value Locked rather than transaction volume, so the corresponding
metric is the \textit{Growth-Levelized NV/TVL ratio}: market capitalization over the
growth-adjusted present value of TVL, built with the same discounting machinery as the
coin metric. Crucially, neither
valuation ratio carries its own independent theoretical prediction for returns. Rather,
each serves as a conditioning variable: the conviction signal is most valuable as a
return predictor when the market has not yet priced it in---that is, when the valuation
ratio is low. Assets that combine high conviction with cheap valuation (Stars) offer the
best expected returns; assets that combine low conviction with expensive valuation
(Avoid) offer the worst. We find that the conviction signal alone positively predicts
future returns cross-sectionally for both coins and governance tokens. The predictive
power is substantially stronger for assets with below-median valuation ratios, consistent
with the view that the valuation ratio identifies when the conviction signal remains
unpriced. A long Stars/short Avoid calendar-time portfolio generates an annualized alpha
of [X]\% with a Sharpe ratio of [Y] after controlling for market, size, and momentum
factors, with the results holding separately for coins and governance tokens.

\paragraph{Related Literature.}
A growing body of work attempts to value digital assets, and our paper connects to
several strands. On equilibrium cryptocurrency valuation, \citet{pagnotta2022equilibrium}
model Bitcoin value as a function of user adoption and network security;
\citet{biais2019blockchain} study equilibrium selection in blockchain protocols;
\citet{schilling2019bitcoin} develop a general equilibrium model of Bitcoin as a competing
currency; and \citet{easley2019mining} analyze how transaction fees connect mining
incentives to market value. \citet{cong2021tokenomics} and \citet{sockin2020model} are
the closest to our approach in modeling token prices through adoption dynamics and
convenience yields respectively; we complement them by providing an empirically tractable
QTM decomposition and a testable investment signal derived from on-chain observables.
On crypto risk factors and return predictability, \citet{liu2021risks} identify
network, momentum, and investor attention factors; \citet{liu2022common} develop a
three-factor model for cryptocurrency returns; and \citet{makarov2020trading} document
price dislocations across exchanges. We contribute to this literature the first
on-chain supply-side signal---$\lambda/(1-\lambda)$---grounded in monetary theory rather
than return-based factor construction. On the monetary economics of digital assets,
\citet{kiyotaki1993search}, \citet{lagos2005unified}, and \citet{schilling2019bitcoin}
provide search-theoretic and general equilibrium foundations; \citet{cong2024staking}
analyzes proof-of-stake tokenomics; and \citet{jiang2024exorbitant} draw the analogy
between dollar exorbitant privilege and dominant-network premia. The convenience yield
literature \citep{krishnamurthy2012aggregate, nagel2016liquidity} motivates the monetary
premium in $\text{MC}_{\text{SoV}}$. On token issuance and governance,
\citet{howell2020initial} study the determinants of ICO success;
\citet{chod2021theory} model ICO design under information asymmetry; and
\citet{edmans2011governance} and \citet{hirschman1970exit} provide the governance
theory linking holder participation to firm and protocol value. Our
paper is the first, to our knowledge, to derive a directly observable valuation
decomposition from QTM, to distinguish the signal content of $\lambda$ for coins versus
governance tokens, and to demonstrate that the interaction of conviction and valuation
predicts cross-sectional cryptocurrency returns.

The remainder of the paper is organized as follows. Section~\ref{sec:theory} develops
the theoretical frameworks for coins (SoV/MoE) and tokens (governance moat) and derives
the testable propositions. Section~\ref{sec:measures} defines and constructs the three
key variables: the conviction index $\lambda$, the Growth-Levelized NVT ratio for coins,
and the Growth-Levelized NV/TVL ratio for tokens. Section~\ref{sec:data} describes the data sources,
sample construction, and summary statistics. Section~\ref{sec:results} presents the
cross-sectional return predictability and portfolio results. Section~\ref{sec:robust}
reports robustness tests. Section~\ref{sec:conclusion} concludes.

%=============================================================
\section{Theoretical Framework}
\label{sec:theory}
%=============================================================

This section develops the theoretical argument in three steps. Section~\ref{sec:theory_coins}
establishes the coin framework: QTM applied to native blockchain coins yields a SoV/MoE
decomposition in which the locking ratio $\lambda$ is the sufficient statistic.
Section~\ref{sec:theory_tokens} develops the token framework: governance tokens lack a
medium-of-exchange function, so $\lambda$ is reinterpreted as a governance moat signal
grounded in costly signaling theory. Section~\ref{sec:theory_valuation} introduces the
common valuation-conditioning mechanism that applies across both tracks. Each subsection
closes with the empirical propositions we test.

\subsection{The Coin Framework: SoV/MoE and the Staking Ratio}
\label{sec:theory_coins}

The QTM identity applied to a native blockchain coin treats market capitalization as
money supply: $\text{MC} \cdot V = PQ$, where $V$ is the velocity of supply in
transactional use and $PQ$ is nominal on-chain output. The critical observation is that
supply is heterogeneous in velocity. A fraction $\lambda$ of total circulating supply is
locked in proof-of-stake validation contracts, staked to earn seigniorage, or otherwise
immobilized with near-zero velocity. The complementary fraction $(1-\lambda)$ is freely
circulating and economically active. Only the active fraction supports transactions;
the locked fraction is held as a store of value and does not circulate. Partitioning
supply accordingly, the appendix shows that market capitalization decomposes as
\[
\text{MC} = \underbrace{\frac{PQ}{V_{\text{active}}}}_{\text{MC}_{\text{MoE}}} + \underbrace{\frac{PQ}{V_{\text{active}}} \cdot \frac{\lambda}{1-\lambda}}_{\text{MC}_{\text{SoV}}},
\]
and therefore the SoV/MoE ratio satisfies
\[
\frac{\text{MC}_{\text{SoV}}}{\text{MC}_{\text{MoE}}} = \frac{\lambda}{1-\lambda}.
\]
This is an accounting identity given $\lambda$, not yet a behavioral prediction.

The behavioral prediction follows from modeling the individual staking decision. Consider
a continuum of agents indexed by $i \in [0,1]$, each choosing at date $t$ whether to
stake one unit of the native coin or hold it freely circulating. Staking yields a publicly
observed current reward $b_t > 0$---seigniorage paid to validators in a PoS network
\citep{cong2024staking}---plus a private convenience yield $c_{i,t+1} = \theta_{t+1} +
\varepsilon_i$, where $\theta_{t+1}$ is the network's true future growth trajectory
(unobserved at $t$) and $\varepsilon_i$ is mean-zero idiosyncratic signal noise.
Staking also imposes a private liquidity cost $\varphi_i$ (i.i.d., independent of
$\theta_{t+1}$ and $\varepsilon_i$). Agent $i$ stakes if and only if
\[
b_t + \theta_{t+1} + \varepsilon_i \geq \varphi_i \quad \Longleftrightarrow \quad
\eta_i \leq b_t + \theta_{t+1}, \qquad \eta_i \equiv \varphi_i - \varepsilon_i.
\]
By the law of large numbers, the aggregate staking ratio converges almost surely to
\[
\lambda_t = H(b_t + \theta_{t+1}),
\]
where $H$ is the CDF of $\eta_i$. Idiosyncratic cost and signal noise wash out in the
cross-section; what survives is a deterministic, strictly increasing function of
$\theta_{t+1}$. Because $b_t$ is publicly observed, $\lambda_t$ is a publicly verifiable
signal of the holder base's aggregate private conviction about the network's future growth.
Extracting $\theta_{t+1}$ from $\lambda_t$, however, requires knowing $H$ and $b_t$ with
precision---not common knowledge for most market participants. As in
\citet{grossman1980impossibility}, the signal is not fully impounded into price at $t$.
As $\theta_{t+1}$ materializes and becomes public over $t+1$, price converges toward the
value anticipated by high-$\lambda$ stakers, generating positive return predictability.

\textbf{Proposition 1a (Coins).} $\mathrm{Cov}_j(r_{j,t+1},\,\lambda_{j,t}/(1-\lambda_{j,t})) > 0$
in the cross-section: coins with higher SoV/MoE ratios earn higher subsequent returns,
after controlling for size, momentum, and market beta.

\textbf{Hypothesis 1a.} \textit{For native proof-of-stake coins, assets with higher
$\lambda/(1-\lambda)$ in month $t$ earn statistically and economically significant
higher returns over month $t+1$, after controlling for coin-level size, momentum, and
market exposure. The result holds regardless of whether the coin is an L1 or L2, provided
the staking mechanism confers a positive seigniorage benefit $b_t > 0$.}

\subsection{The Token Framework: Governance Conviction as Economic Moat}
\label{sec:theory_tokens}

For governance and DeFi protocol tokens, the SoV/MoE framework is inapplicable: these
tokens are not media of exchange. A holder of UNI does not need to spend UNI to use
Uniswap; holding governance tokens does not reduce the protocol's transactional capacity.
The token has no meaningful velocity in the QTM sense, and the SoV/MoE decomposition
carries no economic content. What $\lambda$ signals for tokens is something different:
\textit{governance conviction}---the fraction of supply whose holders have voluntarily
incurred a cost to demonstrate fundamental belief in the protocol's future.

The formal model is a straightforward extension of the coin model. The same continuum of
agents now hold a governance token at date $t$ and choose whether to lock it into a
governance contract (vote-escrow, delegation, or protocol staking) or hold it freely.
Locking carries a private cost $\varphi_i = \ell + \kappa_i$, decomposed into a
common illiquidity cost $\ell \geq 0$ (positive for vote-escrow protocols such as
Curve's veCRV or Balancer's veBAL, zero for simple delegation under ERC20Votes) and an
idiosyncratic engagement cost $\kappa_i \geq 0$, the monitoring and time cost of
exercising governance voice rather than exit \citep{hirschman1970exit, edmans2011governance}.
The current locking benefit $b_t \geq 0$ may include a fee or revenue share for protocols
with on-chain distribution; for most governance tokens $b_t = 0$ and the incentive to
lock is purely forward-looking. The threshold condition is identical to the coin case:
$\eta_i = \varphi_i - \varepsilon_i \leq b_t + \theta_{t+1}$, yielding
$\lambda_t = H(b_t + \theta_{t+1})$.

The key difference is interpretation. For coins, $b_t > 0$ is paid in newly minted
seigniorage, so staking is partly a rational income decision. For tokens with $b_t = 0$,
locking is a purely conviction-revealing action: the holder incurs $\ell + \kappa_i$ with
no current compensation, revealing that their private belief about $\theta_{t+1}$ is strong
enough to make the commitment worthwhile. This is a Spence \citeyearpar{spence1973job}
costly-signaling equilibrium: only holders with sufficiently high conviction rationally pay
the cost. The resulting aggregate $\lambda_t$ is therefore a cleaner signal of fundamental
conviction for governance tokens than for coins, because the seigniorage confound is
absent. The partial-revelation logic applies identically: the signal in $\lambda_t$ is not
freely extractable from public data, so it is not instantly priced, generating the same
return predictability.

\textbf{Proposition 1b (Tokens).} $\mathrm{Cov}_j(r_{j,t+1},\,\lambda_{j,t}) > 0$
in the cross-section of governance tokens: assets with higher governance conviction ratios
earn higher subsequent returns, after controlling for size, momentum, and protocol category.

\textbf{Hypothesis 1b.} \textit{For DeFi governance tokens, assets with higher $\lambda$
in month $t$ earn statistically and economically significant higher returns over month
$t+1$, after controlling for token-level size, momentum, and protocol category. The
predictive power of voting-participation-weighted $\lambda$ exceeds that of passive
locking alone, consistent with the costly-signaling logic: the harder the action, the
purer the signal.}

\subsection{The Valuation Conditioning Signal}
\label{sec:theory_valuation}

Propositions 1a and 1b establish that $\lambda$ carries information not yet fully in price.
What they do not characterize is how much of the signal remains unpriced at a given date.
Let $\delta_t \in [0,1]$ denote the fraction of the conviction signal in $\lambda_t$
already absorbed into the current price, so that $p_t = p_t^{\text{fund}}(1+\delta_t)$
where $p_t^{\text{fund}}$ reflects current economic activity. The expected return as the
signal is validated over $t+1$ is proportional to the unabsorbed portion
$\lambda_t(1 - \delta_t)$, not to $\lambda_t$ alone. The marginal predictive power of
$\lambda_t$ is therefore decreasing in $\delta_t$, and $\delta_t$ is increasing in the
current market capitalization relative to fundamental activity.

For coins, the relevant fundamental is discounted expected on-chain volume: $\delta_t$
is increasing in $\text{NVT\_GL}_t = \text{MC}_t / PQ_t^*$, where $PQ^*_t$ is the
growth-adjusted present value of expected transaction activity (formally defined in
Section~\ref{sec:measures_nvt}). For tokens, the relevant fundamental is the
growth-adjusted present value of Total Value Locked: $\delta_t$ is increasing in
$\text{NV/TVL\_GL}_t = \text{MC}_t / \text{TVL}^*_t$, where $\text{TVL}^*_t$ is
constructed with the same discounting machinery (Section~\ref{sec:measures_tvl}).
In both cases, a high ratio indicates that current market capitalization already embeds
the conviction signal, and the incremental return predictability of $\lambda_t$ is low;
a low ratio indicates the signal remains unabsorbed.

\textbf{Proposition 2.} \textit{The cross-partial effect of $\lambda_t$ and the
valuation ratio on expected returns is negative: the return predictability of $\lambda_t$
is concentrated among low-valuation-ratio assets.}

\textbf{Hypothesis 2.} \textit{The return predictability documented in H1a and H1b is
significantly stronger among assets with below-median valuation ratios (NVT\_GL for coins,
NV/TVL\_GL for tokens). The marginal return to a unit increase in $\lambda$ is higher when
the asset is cheap relative to its fundamental economic activity.}

The valuation ratio carries no independent return prediction in this framework. A low
NVT\_GL alone---without a high $\lambda$---merely describes a coin with limited
current throughput relative to market cap, not necessarily a mispricing; similarly for
NV/TVL\_GL. The conditioning role is asymmetric: valuation identifies the \textit{timing}
of conviction signal absorption, not a separate source of alpha.

\subsection{Hypotheses and the Quadrant Portfolio}
\label{sec:theory_portfolio}

Propositions 1 and 2 jointly define a two-dimensional framework. \textit{Stars} combine
high $\lambda$ with low valuation ratio---strong conviction, signal not yet priced---and
are predicted to have the highest future returns. \textit{Avoid} assets combine low
$\lambda$ with high valuation ratio---weak conviction, valuation already stretched---and
are predicted to have the lowest returns.

\textbf{Hypothesis 3.} \textit{A long-short portfolio that goes long Star assets
(above-median $\lambda$, below-median valuation ratio) and short Avoid assets
(below-median $\lambda$, above-median valuation ratio), rebalanced monthly, generates
positive and statistically significant factor-adjusted alpha after controlling for market,
size, and momentum. The result holds separately within the coin sample (using NVT\_GL)
and within the token sample (using NV/TVL\_GL), and in the combined panel. The Stars-minus-Avoid
Sharpe ratio exceeds that of equal-weighted benchmark portfolios.}

H3 is the integrating test: it asks whether the joint two-dimensional signal has
incremental predictive power beyond either dimension alone. We report results separately
for coins and tokens to preserve the track distinction: for coins, alpha reflects the
return to SoV/MoE mispricing correcting over time; for tokens, alpha reflects the market's
underreaction to governance conviction signals.

%=============================================================
\section{Variable Construction}
\label{sec:measures}
%=============================================================

The theoretical framework in Section~\ref{sec:theory} requires three empirical proxies:
the conviction index $\lambda$, common to both asset tracks; the Growth-Levelized NVT
ratio for the coin track; and the Growth-Levelized NV/TVL ratio for the token track. This section defines
each variable and describes its construction. Data sources and sample formation are
covered in Section~\ref{sec:data}.

\subsection{The Conviction Index \texorpdfstring{$\lambda$}{λ}}
\label{sec:measures_lambda}

The theoretical $\lambda_t$ is the aggregate locking or commitment ratio of an asset's
holder base at date $t$. Empirically, the conviction-revealing actions observable on-chain
vary by asset type and protocol design. We therefore construct $\lambda$ as a composite
index over three conceptually motivated channels, and aggregate across available channels
using cross-sectional z-scores.

\paragraph{Channel 1: Staking.}
For proof-of-stake layer-one coins, the staking channel is the fraction of circulating
supply bonded as validator or delegator stake at month-end. This is the direct empirical
counterpart of the locking action in Section~\ref{sec:theory_coins}: it simultaneously
earns seigniorage ($b_t > 0$) and constitutes a costly, publicly visible commitment to
the network's future. Staking rates are sourced from chain-native archive data for each
coin.\footnote{For Polkadot (DOT) and Kusama (KSM), bonded supply is retrieved from the
Subscan archive API. For Cosmos-ecosystem chains (CRO, KAVA, OSMO), we query the Cosmos
SDK Staking module endpoint at block heights corresponding to month-end timestamps, with
timestamps identified via binary search on the CometBFT RPC interface. For EVM-based
staking contracts (BNB, MATIC, TRX, and others), staked supply is derived from on-chain
staking event logs via the Etherscan multi-chain archive API. Each data source is
validated by confirming that reported values vary across block heights, ensuring the
archive reflects historical rather than current-only state.} The staking channel covers
48 assets and 2,222 asset-months.

\paragraph{Channel 2: Holding.}
The holding channel measures the fraction of circulating supply that has not transacted
on-chain in the preceding six months (HODL-6m). This proxy extends conviction measurement
to any EVM-compatible asset without a formal staking mechanism, mapping to the inert
supply fraction $\lambda_t$ in the model: a holder who has not transacted for six months
has revealed a preference for holding over liquidity. HODL-6m is computed by replaying
on-chain Transfer event logs in chronological order under a FIFO attribution rule,
assigning each unit of supply to its most recent transfer date.\footnote{Transfer events
are retrieved from the Etherscan Pro V2 multi-chain archive, covering Ethereum mainnet and
twelve additional EVM chains. Two data quality filters are applied: (i) monthly on-chain
flow exceeding 100$\times$ reported circulating supply is classified as a minting,
bridging, or aggregator artifact and excluded; (ii) asset-months in which HODL-6m exceeds
80\% persistently are flagged for robustness checks, consistent with possible custodial
cold-wallet concentration rather than individual holder conviction. Self-transfers (sender
equals recipient) are excluded from the event replay.} This channel spans 408 assets and
11,654 asset-months, providing the broadest coverage in the conviction panel.

\paragraph{Channel 3: Governance.}
Two governance sub-channels capture active participation in protocol decision-making.
The delegation sub-channel records the fraction of circulating supply delegated to a
voting representative under an ERC20Votes contract, observable from on-chain Delegate
event logs.\footnote{Data sourced from \texttt{DelegateVotesChanged} event logs via the
Etherscan Pro V2 API, aggregated to month-end totals.} The voting sub-channel records
proposal participation rates from Snapshot off-chain governance records, cross-referenced
with on-chain ERC20Votes logs where available. Both sub-channels identify holders who
have incurred the engagement costs $\kappa_i$ of Section~\ref{sec:theory_tokens}---the
costlier the action, the more selected and informative the resulting signal. The
delegation and voting sub-channels cover 26 and 53 assets, respectively, concentrated in
the DeFi governance token sample.

\paragraph{Composite Index.}
For each asset-month, each available channel is standardized to a z-score over all
asset-months with non-missing observations for that channel. The composite $\lambda_z$ is
the equal-weight average of standardized channel scores available for that asset-month.
The result has approximately zero cross-sectional mean and is defined whenever at least
one channel is observed. Equal weighting imposes no prior on the relative information
content of the channels; channel-specific predictive power is examined in robustness
tests. The full conviction panel covers 467 assets and 13,626 asset-months, of which
11,901 are single-channel, 1,371 two-channel, and 354 three-or-more-channel observations.

\subsection{Growth-Levelized NVT}
\label{sec:measures_nvt}

For the coin regression track, the valuation anchor is the Growth-Levelized Network Value
to Transactions ratio (NVT\_GL), which scales market capitalization by the
growth-adjusted, risk-discounted present value of expected on-chain throughput:
\begin{align}
    \text{NVT\_GL}_{j,t} &= \frac{\text{MC}_{j,t}}{PQ^*_{j,t}}, \label{eq:nvt_gl} \\[6pt]
    PQ^*_{j,t} &= \frac{\displaystyle
        \sum_{s=1}^{n} \frac{PQ_{0,j,t}(1+g_{j,t})^{\!s}}{(1+r_{e,j,t})^{\!s}}
        + \frac{PQ_{0,j,t}(1+g_{j,t})^{\!n}(1+g_\infty)}{(r_{e,j,t}-g_\infty)(1+r_{e,j,t})^{\!n}}
    }{\displaystyle
        \frac{1-(1+r_{e,j,t})^{-n}}{r_{e,j,t}}
    }, \label{eq:pq_star}
\end{align}
where $PQ_{0,j,t}$ is the trailing twelve-month annualized on-chain transfer volume for
asset $j$ at month $t$; $g_{j,t}$ is the trailing three-year compound annual growth rate
of $PQ_0$, capped in $[-50\%,\,200\%]$; $n = 10$ years is the explicit-growth projection
horizon; $g_\infty = 3\%$ is the terminal perpetuity growth rate; and
$r_{e,j,t} = r_f + \hat{\beta}_{j,t} \cdot \text{MRP}$ is the asset-specific discount
rate, with $r_f = 4\%$, $\text{MRP} = 30\%$, and $\hat{\beta}_{j,t}$ the trailing
36-month beta against the BTC price index, floored so that $r_{e,j,t} \geq g_\infty + 2\%$
to ensure well-defined terminal values.\footnote{The denominator of \eqref{eq:pq_star}
is the standard annuity factor normalizing the finite-horizon present value of future $PQ$
to an equivalent annualized flow, making NVT\_GL interpretable as a
price-to-normalized-activity ratio analogous to a price-to-normalized-earnings multiple
in equity analysis. The MRP of 30\% reflects the elevated required return on
cryptocurrency risk; $\hat\beta_{j,t}$ is retained in the panel so alternative discount
rate assumptions can be applied directly in robustness tests.}

NVT\_GL below one signals that market capitalization is below the growth-adjusted present
value of expected on-chain activity. NVT\_GL substantially above one indicates that
current market capitalization already embeds forward growth not yet visible in throughput.
The growth adjustment corrects a systematic bias in raw NVT: fast-growing assets may
appear expensive on a raw NVT basis but inexpensive once expected growth is discounted.

\subsection{Growth-Levelized NV/TVL}
\label{sec:measures_tvl}

For the token regression track, the valuation anchor is the Growth-Levelized Network
Value to Total Value Locked ratio, constructed with the same discounted-present-value
machinery as NVT\_GL. Total Value Locked---the market value of assets deposited in a
protocol's smart contracts---proxies for the scale of economic activity the protocol
intermediates, and is the token analog of on-chain transaction volume in the QTM
framework. Because TVL is a stock rather than a flow, the base is a trailing average
level rather than a trailing sum:
\begin{align}
    \text{NV/TVL\_GL}_{j,t} &= \frac{\text{MC}_{j,t}}{\text{TVL}^*_{j,t}}, \label{eq:nv_tvl_gl} \\[6pt]
    \text{TVL}^*_{j,t} &= \frac{\displaystyle
        \sum_{s=1}^{n} \frac{\text{TVL}_{0,j,t}(1+g_{j,t})^{\!s}}{(1+r_{e,j,t})^{\!s}}
        + \frac{\text{TVL}_{0,j,t}(1+g_{j,t})^{\!n}(1+g_\infty)}{(r_{e,j,t}-g_\infty)(1+r_{e,j,t})^{\!n}}
    }{\displaystyle
        \frac{1-(1+r_{e,j,t})^{-n}}{r_{e,j,t}}
    }, \label{eq:tvl_star}
\end{align}
where $\text{TVL}_{0,j,t}$ is the trailing twelve-month average of month-end TVL for
protocol $j$ (requiring at least six monthly observations); $g_{j,t}$ is the trailing
three-year compound annual growth rate of $\text{TVL}_0$, with two- and one-year
fallbacks and the same $[-50\%,\,200\%]$ cap; and all remaining
parameters---$n$, $g_\infty$, $r_f$, MRP, and the beta construction---are identical to
those in \eqref{eq:pq_star}, so the coin and token tracks share a single set of
valuation assumptions. Where a protocol is deployed across multiple chains, TVL is the
sum of all-chain TVL.

The growth adjustment serves the same purpose as in the coin track: a protocol whose TVL
is expanding rapidly may appear expensive on a raw NV/TVL basis but inexpensive once
expected growth is priced in, and conversely a protocol with collapsing TVL may appear
deceptively cheap on the raw ratio. A low NV/TVL\_GL indicates that market capitalization
is modest relative to the growth-adjusted present value of the protocol's economic
footprint; a high NV/TVL\_GL indicates the reverse.

%=============================================================
\section{Data}
\label{sec:data}
%=============================================================

\subsection{Data Sources}
\label{sec:data_sources}

Price, market capitalization, and circulating supply are drawn from CoinMarketCap at
daily granularity, snapped to end-of-calendar-month observations. CoinMarketCap provides
the broadest historical coverage of digital asset market data and is the standard
reference for universe construction in the cryptocurrency finance literature
\citep{liu2022common,makarov2020trading}.

On-chain Transfer event logs, staking event logs, and ERC20Votes Delegate event logs
are retrieved from the Etherscan Pro V2 multi-chain API, which provides archive access
to EVM-compatible chains including Ethereum mainnet, BNB Chain, Polygon, Avalanche,
Arbitrum, Optimism, and nine additional networks. Non-EVM staking data follows
chain-native sources: Subscan archive API for Polkadot and Kusama; Cosmos SDK Staking
module endpoints for Cosmos-ecosystem chains. Off-chain governance participation is
sourced from the Snapshot API, cross-referenced with on-chain ERC20Votes logs.

On-chain transaction volume for the NVT\_GL construction is assembled from three sources
in priority order. DeFiLlama chain-level DEX
volume\footnote{\url{https://defillama.com}, accessed through May 2026.}
serves as the primary source for chains with active decentralized exchange infrastructure
(1,340 asset-months). BitInfoCharts sent-in-USD settlement
value\footnote{\url{https://bitinfocharts.com}, accessed through May 2026.}
is used for PoW coins and pre-DEX periods, capturing all on-chain transfers rather than
DEX activity alone (705 asset-months). Blockchain.com estimated transaction
value\footnote{\url{https://blockchain.com}, accessed through May 2026.}
supplements Bitcoin and legacy chain observations with pre-DEX-era coverage (118
asset-months). The remaining 363 asset-months use DeFiLlama protocol-level volume for
assets whose primary activity flows through specific protocols. Each series is validated
against the chain's own reported aggregate, with the reconstructed monthly level required
to lie within 5\% of the chain total. The resulting NVT\_GL panel covers 67 assets and
2,526 valid asset-months from August 2016 through May 2026.

Total Value Locked is sourced from DeFiLlama at daily granularity, snapped to month-end;
multi-chain protocols use the sum of all-chain TVL. The raw TVL panel covers 162 assets
and 8,066 asset-months from September 2017 through May 2026. Applying the
growth-levelization of \eqref{eq:tvl_star}---which requires a trailing twelve-month TVL
average and at least one year of prior history for the growth estimate---yields an
NV/TVL\_GL panel of 111 assets and 3,125 asset-months from January 2021 through May 2026.

\subsection{Sample Construction}
\label{sec:data_sample}

The starting universe is 1,939 digital assets drawn from the CoinMarketCap historical
database, observed at end-of-calendar-month from August 2015 through May 2026, yielding
156,838 asset-month observations.\footnote{All assets with at least one non-zero price
observation in CoinMarketCap are included. Market capitalization is computed as reported
price times circulating supply. Continuous coverage is not required; assets may enter and
exit the panel as they gain or lose listing.}

Assets are classified into three mutually exclusive types: \textit{coins} (native
settlement assets of a public blockchain, e.g., BTC, ETH, SOL), \textit{tokens}
(smart-contract-issued assets deployed on a host chain, e.g., ERC-20 governance and DeFi
tokens), and \textit{other} (stablecoins, wrapped tokens, and infrastructure assets whose
value is structurally indexed to an external reference or operating entity).
Other-class assets are excluded from both regression samples because neither the SoV/MoE
framework nor the governance moat framework applies to their valuation.

The two regression samples are defined by simultaneous availability of $\lambda_z$ and
a positive valuation denominator. The \textit{coin regression sample} requires: (i) coin
classification; (ii) proof-of-stake architecture; (iii) at least one month of non-missing
$\lambda_z$; and (iv) at least one month of positive NVT\_GL. This yields 24 assets and
718 asset-months spanning December 2020 through May 2026. The median NVT\_GL in this
sample is 0.017 and median market capitalization is \$820 million. The \textit{token
regression sample} requires: (i) token classification; (ii) non-missing $\lambda_z$; and
(iii) positive NV/TVL\_GL. This yields 101 assets and 2,771 asset-months spanning
January 2021 through May 2026. Median NV/TVL\_GL is 2.29 (IQR: 0.20--16.10), median TVL
is \$42.2 million, and median market capitalization is \$98 million.\footnote{The
growth-levelization requirement shortens the token sample relative to a raw NV/TVL
alternative (136 assets, 4,627 asset-months from March 2020): assets without at least one
year of prior TVL history have no growth estimate and drop out. The raw NV/TVL sample is
retained for robustness comparisons.}

Several populations are excluded from one or both samples by data structure. Proof-of-work
coins have no staking mechanism and are therefore outside the scope of Hypothesis~1a,
though they contribute to NVT\_GL baseline tests. Among PoS coins, historical staking
data is unavailable for chains whose archive nodes do not retain historical state, either
due to aggressive pruning or reliance on third-party indexers for pre-recent-block queries.
Among tokens, approximately 1,100 assets have no DeFiLlama TVL entry because they serve as
utility, memetic, or infrastructure tokens for which protocol TVL is not defined. The
remaining coverage gap arises from temporal non-overlap: tokens for which conviction data
is available only after the period of active TVL reporting contribute to individual panels
but not to the regression sample intersection. Channel 2 (holding) applies only to
EVM-compatible assets; PoW coins with UTXO-based ledger accounting require a different
pipeline based on UTXO age analysis and are left to future work.

The combined regression panel spans 125 unique assets and 3,489 asset-months. At the end
of the sample in May 2026, the panel accounts for \$451 billion in market
capitalization---19.0\% of the \$2.38 trillion CoinMarketCap universe---consistent with
the concentration of on-chain governance and staking infrastructure among the larger
assets in the market.

\subsection{Summary Statistics}
\label{sec:data_stats}

Table~\ref{tab:summary} reports coverage across data panels and key statistics for the
regression samples.

\begin{table}[ht]
\centering
\caption{Summary Statistics}
\label{tab:summary}
\small
\begin{tabular}{lrrcc}
\toprule
Panel / Sample & Assets & Asset-months & Date range & Median MC \\
\midrule
Full universe              & 1,939 & 156,838 & 2015-08 to 2026-05 & \$6.5M \\
Conviction panel ($\lambda_z$) & 467   & 13,626  & 2016-12 to 2026-05 & --- \\
\quad ch1 staking          & 48    & 2,222   & 2017-07 to 2026-05 & --- \\
\quad ch2 holding          & 408   & 11,654  & 2015-08 to 2026-05 & --- \\
\quad ch3 delegation       & 26    & 627     & 2020-05 to 2026-05 & --- \\
\quad ch3 voting           & 53    & 1,202   & 2020-08 to 2026-05 & --- \\
NVT\_GL panel              & 67    & 2,526   & 2016-08 to 2026-05 & --- \\
TVL panel (raw)            & 162   & 8,066   & 2017-09 to 2026-05 & --- \\
NV/TVL\_GL panel           & 111   & 3,125   & 2021-01 to 2026-05 & --- \\
\midrule
Coin regression (NVT\_GL)  & 24    & 718     & 2020-12 to 2026-05 & \$820M \\
Token regression (NV/TVL\_GL) & 101 & 2,771  & 2021-01 to 2026-05 & \$98M \\
Combined                   & 125   & 3,489   & 2020-12 to 2026-05 & \$138M \\
\bottomrule
\end{tabular}
\smallskip

\parbox{0.9\linewidth}{\footnotesize
\textit{Notes:} Median MC is the cross-sectional median market capitalization across all
asset-months with a positive price observation. Conviction panel sub-channel counts sum
to more than the panel total because 1,725 asset-months appear in two or more channels.
Coin regression requires simultaneous non-missing $\lambda_z$ and positive NVT\_GL for
proof-of-stake coins. Token regression requires simultaneous non-missing $\lambda_z$ and
positive NV/TVL\_GL. Median NVT\_GL in the coin regression sample is 0.017. Median
NV/TVL\_GL in the token regression sample is 2.29 (IQR: 0.20--16.10); median TVL is
\$42.2M.
}
\end{table}

%=============================================================
\section{Results}
\label{sec:results}
%=============================================================

\textit{[To be written.]}

%=============================================================
\section{Robustness}
\label{sec:robust}
%=============================================================

\textit{[To be written.]}

%=============================================================
\section{Conclusion}
\label{sec:conclusion}
%=============================================================

\textit{[To be written.]}

%=============================================================
\appendix
\section{Derivation of the SoV/MoE Decomposition}
\label{app:derivations}
%=============================================================

This appendix provides the formal derivation of the SoV/MoE decomposition and its
simplification to a function of $\lambda$ alone.

\paragraph{Setup.} Let $M$ denote the total circulating supply of a digital asset and
$\text{MC}$ its market capitalization, which serves as the money supply proxy in the
QTM framework. Let $PQ$ denote nominal economic output (on-chain transaction volume for
coins; protocol throughput for tokens). Standard QTM gives $\text{MC} \cdot V = PQ$,
so the market capitalization satisfies $\text{MC} = PQ / V$.

\paragraph{The locking ratio.} Define $\lambda$ as the proportion of total supply that
is locked or staked:
\[
\lambda = \frac{M_{\text{locked}}}{M}, \qquad \lambda \in [0,1).
\]
The free-circulating supply is $M(1-\lambda)$. We define the \textit{free velocity}
$V_{\text{free}}$ as the velocity of only the economically active (unlocked) portion:
\[
V_{\text{free}} = \frac{PQ}{M(1-\lambda)}.
\]

\paragraph{MoE and SoV components.} The Medium of Exchange component of market
capitalization is the value attributable to current transactional activity, priced at
the free velocity:
\[
\text{MC}_{\text{MoE}} = \frac{PQ}{V_{\text{free}}} = M(1-\lambda).
\]
The Store of Value component is the residual:
\[
\text{MC}_{\text{SoV}} = \text{MC} - \text{MC}_{\text{MoE}} = M - M(1-\lambda) = M\lambda.
\]

\paragraph{The SoV/MoE ratio.} The ratio of the SoV to MoE components is:
\[
\frac{\text{MC}_{\text{SoV}}}{\text{MC}_{\text{MoE}}} = \frac{M\lambda}{M(1-\lambda)}
= \frac{\lambda}{1-\lambda}.
\]
This result shows that the SoV/MoE ratio is a monotone increasing function of $\lambda$
alone, requiring no estimation of velocity, transaction volume, or any other derived
quantity. As $\lambda \to 0$, the ratio approaches zero (pure medium of exchange). As
$\lambda \to 1$, the ratio diverges (pure store of value).

\paragraph{The Growth-Levelized NVT ratio.} The Network Value to Transactions ratio is
$\text{NVT} = \text{MC} / PQ_{\text{annualized}}$. We construct a forward-looking
version by replacing current $PQ$ with $PQ^*$, the levelized present value of the
projected transaction stream:
\[
PQ^* = \frac{\sum_{s=1}^{n} \frac{PQ_0 (1+g)^s}{(1+r_e)^s}
+ \frac{PQ_0(1+g)^n(1+g_{\infty})}{(r_e - g_{\infty})(1+r_e)^n}}
{\text{annuity factor}(r_e, n)},
\]
where $g$ is the trailing $k$-year CAGR of $PQ$, $r_e$ is the CAPM-estimated required
return, $g_{\infty}$ is the terminal growth rate, and the annuity factor converts the
present value to a levelized annual equivalent (analogous to the Excel PMT function).
The Growth-Levelized NVT ratio is then $\text{NVT}_{GL} = \text{MC} / PQ^*$. Full estimation
details are in Section~\ref{sec:data}.

%=============================================================
%  BIBLIOGRAPHY
%=============================================================
\bibliographystyle{jf}
% \bibliography{references}

\begin{thebibliography}{99}

\bibitem[BitInfoCharts(2026)]{bitinfocharts2026}
BitInfoCharts, 2026, Blockchain statistics and sent-value data,
\url{https://bitinfocharts.com}, accessed May 2026.

\bibitem[Blockchain.com(2026)]{blockchain2026}
Blockchain.com, 2026, Bitcoin blockchain explorer and transaction statistics,
\url{https://blockchain.com}, accessed May 2026.

\bibitem[Biais, Bisière, Bouvard, and Casamatta(2019)]{biais2019blockchain}
Biais, B., C.~Bisière, M.~Bouvard, and C.~Casamatta, 2019, The blockchain folk theorem,
\textit{Review of Economic Studies} 86, 1341--1371.

\bibitem[CoinMarketCap(2026)]{cmc2026}
CoinMarketCap, 2026, Historical cryptocurrency market data,
\url{https://coinmarketcap.com}, accessed May 2026.

\bibitem[Chod and Lyandres(2021)]{chod2021theory}
Chod, J., and E.~Lyandres, 2021, A theory of ICOs: Diversification, agency, and
information asymmetry, \textit{Management Science} 67, 5969--5989.

\bibitem[Cong, Li, and Wang(2021)]{cong2021tokenomics}
Cong, L.~W., Y.~Li, and N.~Wang, 2021, Tokenomics: Dynamic adoption and valuation,
\textit{Review of Financial Studies} 34, 1105--1155.

\bibitem[Cong and He(2024)]{cong2024staking}
Cong, L.~W., and Z.~He, 2024, The tokenomics of staking,
\textit{NBER Working Paper} 33640.

\bibitem[DeFiLlama(2026)]{defillama2026}
DeFiLlama, 2026, Protocol TVL and decentralized exchange volume data,
\url{https://defillama.com}, accessed May 2026.

\bibitem[Easley, O'Hara, and Basu(2019)]{easley2019mining}
Easley, D., M.~O'Hara, and S.~Basu, 2019, From mining to markets: The evolution of
Bitcoin transaction fees, \textit{Journal of Financial Economics} 134, 91--109.

\bibitem[Edmans and Manso(2011)]{edmans2011governance}
Edmans, A., and G.~Manso, 2011, Governance through exit and voice: A theory of
multiple blockholders, \textit{Review of Financial Studies} 24, 2395--2428.

\bibitem[Fisher(1911)]{fisher1911}
Fisher, I., 1911, \textit{The Purchasing Power of Money},
Macmillan, New York.

\bibitem[Grossman and Stiglitz(1980)]{grossman1980impossibility}
Grossman, S.~J., and J.~E.~Stiglitz, 1980, On the impossibility of informationally
efficient markets, \textit{American Economic Review} 70, 393--408.

\bibitem[Hirschman(1970)]{hirschman1970exit}
Hirschman, A.~O., 1970, \textit{Exit, Voice, and Loyalty},
Harvard University Press, Cambridge, MA.

\bibitem[Howell, Niessner, and Yermack(2020)]{howell2020initial}
Howell, S.~T., M.~Niessner, and D.~Yermack, 2020, Initial coin offerings: Financing
growth with cryptocurrency token sales, \textit{Review of Financial Studies} 33,
3925--3974.

\bibitem[Jiang, Krishnamurthy, and Lustig(2024)]{jiang2024exorbitant}
Jiang, Z., A.~Krishnamurthy, and H.~Lustig, 2024, Exorbitant privilege: A safe-asset
view, \textit{NBER Working Paper} 32454.

\bibitem[Kiyotaki and Wright(1993)]{kiyotaki1993search}
Kiyotaki, N., and R.~Wright, 1993, A search-theoretic approach to monetary economics,
\textit{American Economic Review} 83, 63--77.

\bibitem[Krishnamurthy and Vissing-Jorgensen(2012)]{krishnamurthy2012aggregate}
Krishnamurthy, A., and A.~Vissing-Jorgensen, 2012, The aggregate demand for Treasury
debt, \textit{Journal of Political Economy} 120, 233--267.

\bibitem[Lagos and Wright(2005)]{lagos2005unified}
Lagos, R., and R.~Wright, 2005, A unified framework for monetary theory and policy
analysis, \textit{Journal of Political Economy} 113, 463--484.

\bibitem[Liu and Tsyvinski(2021)]{liu2021risks}
Liu, Y., and A.~Tsyvinski, 2021, Risks and returns of cryptocurrency,
\textit{Review of Financial Studies} 34, 2689--2727.

\bibitem[Liu, Tsyvinski, and Wu(2022)]{liu2022common}
Liu, Y., A.~Tsyvinski, and X.~Wu, 2022, Common risk factors in cryptocurrency,
\textit{Journal of Finance} 77, 1133--1177.

\bibitem[Makarov and Schoar(2020)]{makarov2020trading}
Makarov, I., and A.~Schoar, 2020, Trading and arbitrage in cryptocurrency markets,
\textit{Journal of Financial Economics} 135, 293--319.

\bibitem[Nagel(2016)]{nagel2016liquidity}
Nagel, S., 2016, The liquidity premium of near-money assets,
\textit{Quarterly Journal of Economics} 131, 1927--1971.

\bibitem[Pagnotta and Buraschi(2022)]{pagnotta2022equilibrium}
Pagnotta, E., and A.~Buraschi, 2022, An equilibrium valuation of Bitcoin and
decentralized network assets, \textit{Journal of Finance} 77, 887--943.

\bibitem[Schilling and Uhlig(2019)]{schilling2019bitcoin}
Schilling, L., and H.~Uhlig, 2019, Some simple Bitcoin economics,
\textit{Journal of Monetary Economics} 106, 16--26.

\bibitem[Sockin and Xiong(2020)]{sockin2020model}
Sockin, M., and W.~Xiong, 2020, A model of cryptocurrencies,
\textit{NBER Working Paper} 26816.

\bibitem[Spence(1973)]{spence1973job}
Spence, M., 1973, Job market signaling,
\textit{Quarterly Journal of Economics} 87, 355--374.

\end{thebibliography}

\end{document}
